% Options for packages loaded elsewhere
\PassOptionsToPackage{unicode}{hyperref}
\PassOptionsToPackage{hyphens}{url}
\PassOptionsToPackage{dvipsnames,svgnames,x11names}{xcolor}
\PassOptionsToPackage{space}{xeCJK}
\documentclass[
]{article}
\usepackage{xcolor}
\usepackage[margin=2.5cm]{geometry}
\usepackage{amsmath,amssymb}
\setcounter{secnumdepth}{-\maxdimen} % remove section numbering
\usepackage{iftex}
\ifPDFTeX
  \usepackage[T1]{fontenc}
  \usepackage[utf8]{inputenc}
  \usepackage{textcomp} % provide euro and other symbols
\else % if luatex or xetex
  \usepackage{unicode-math} % this also loads fontspec
  \defaultfontfeatures{Scale=MatchLowercase}
  \defaultfontfeatures[\rmfamily]{Ligatures=TeX,Scale=1}
\fi
\usepackage{lmodern}
\ifPDFTeX\else
  % xetex/luatex font selection
    \setmainfont[]{DejaVu Sans}
    \setmonofont[]{DejaVu Sans Mono}
  \ifXeTeX
    \usepackage{xeCJK}
    \setCJKmainfont[]{Noto Sans CJK SC}
          \fi
  \ifLuaTeX
    \usepackage[]{luatexja-fontspec}
    \setmainjfont[]{Noto Sans CJK SC}
  \fi
\fi
% Use upquote if available, for straight quotes in verbatim environments
\IfFileExists{upquote.sty}{\usepackage{upquote}}{}
\IfFileExists{microtype.sty}{% use microtype if available
  \usepackage[]{microtype}
  \UseMicrotypeSet[protrusion]{basicmath} % disable protrusion for tt fonts
}{}
\makeatletter
\@ifundefined{KOMAClassName}{% if non-KOMA class
  \IfFileExists{parskip.sty}{%
    \usepackage{parskip}
  }{% else
    \setlength{\parindent}{0pt}
    \setlength{\parskip}{6pt plus 2pt minus 1pt}}
}{% if KOMA class
  \KOMAoptions{parskip=half}}
\makeatother
\usepackage{longtable,booktabs,array}
\usepackage{calc} % for calculating minipage widths
% Correct order of tables after \paragraph or \subparagraph
\usepackage{etoolbox}
\makeatletter
\patchcmd\longtable{\par}{\if@noskipsec\mbox{}\fi\par}{}{}
\makeatother
% Allow footnotes in longtable head/foot
\IfFileExists{footnotehyper.sty}{\usepackage{footnotehyper}}{\usepackage{footnote}}
\makesavenoteenv{longtable}
\setlength{\emergencystretch}{3em} % prevent overfull lines
\providecommand{\tightlist}{%
  \setlength{\itemsep}{0pt}\setlength{\parskip}{0pt}}
\usepackage{bookmark}
\IfFileExists{xurl.sty}{\usepackage{xurl}}{} % add URL line breaks if available
\urlstyle{same}
\hypersetup{
  colorlinks=true,
  linkcolor={Maroon},
  filecolor={Maroon},
  citecolor={Blue},
  urlcolor={Blue},
  pdfcreator={LaTeX via pandoc}}

\author{}
\date{}

\begin{document}

\section{直觉与结构 token 的展开序列:
一种语言学解耦}\label{ux76f4ux89c9ux4e0eux7ed3ux6784-token-ux7684ux5c55ux5f00ux5e8fux5217-ux4e00ux79cdux8bedux8a00ux5b66ux89e3ux8026}

\textbf{token 体系为什么能产生可解释的结果, 直觉究竟越过了什么}

\emph{2026-08-12 · 理论提案 (基于 2026-08-12 session 的经验观察,
非已证定理)}

\begin{center}\rule{0.5\linewidth}{0.5pt}\end{center}

\subsection{摘要}\label{ux6458ux8981}

本文提出一个语言学框架, 解耦 token 体系为何能产生可解释结果的
机制与直觉快速路径的机制。核心对比: \textbf{直觉定义 token}(直觉直接
指向的 token)与\textbf{结构 token}(沿定义链展开的 token)的展开序列差异。
论点:

\begin{enumerate}
\def\labelenumi{\arabic{enumi}.}
\tightlist
\item
  token 体系 (B/C/S/G/P 五层: 公设基/概念/符号/语法/呈现) 的可解释
  性来自\textbf{概念层 (C) 与表达层 (S/G/P) 的可分离性} --- 展开序列可以
  在表达层压缩而不改变概念指向, 结果因此可追溯、可解释;
\item
  \textbf{直觉越过的不是概念, 是表达层的展开细节} --- 直觉从概念层直接
  接入, 跳过符号/语法/呈现的生成过程;
\item
  \textbf{直觉跳跃的正确性必须经结构展开验证} (Lean 形式化即此验证) ---
  跳过展开 ≠ 免除展开, 只是延后展开;
\item
  \textbf{核心论点 (普遍联系与捷径)}: 事物的普遍联系 (概念间的连接 网络)
  使直觉能在符号系统下找到连接的捷径, 并且 token 体系允许
  \textbf{直接定义捷径} --- 把直觉发现的捷径固化为可引用的定义 (新 token
  或连接), 使捷径成为结构的一部分而非一次性跳跃;
\item
  \textbf{观测方法 (符号定义捷径的语言学观测)}: 通过对比结构 token 与
  直觉 token 的展开序列, 观测一词多义、多词一义、一义多词等符号
  定义捷径; 核心例: 3+2+7+8 的三种路径 (顺序展开/凑十结构/直觉 直跳) ---
  \textbf{结构与直觉的路径非唯一, 甚至不必须走近路}; 语义理解 =
  在自己的符号-逻辑映射中找到最快路径的数学问题;
\item
  经验证据: 2026-08-12 session (直觉引导的黎曼方向形式化, C011-C025,
  全部 KNOWN 重述, 700k token, 99.2\% 缓存命中, 睡眠后 直觉 compact)。
\end{enumerate}

\subsection{1. 引言}\label{ux5f15ux8a00}

语言认知中一个基本问题: 为什么直觉能''直接知道''而无需完整推导?
本项目观察到: 直觉引导的形式化 (直觉 → Lean 定理) 在 token 经济上 高效,
且直觉陈述总是命中已知结构 (KNOWN 重述)。本文试图解耦这一 现象的机制:
token 体系 (本项目的形式语言) 为何正确, 直觉究竟越过 了什么。

\textbf{记号} (本项目 token 体系): 每个 token 是五层投影的唯一信息元: B
(公设基, 展开终点)、C (概念, 语义核心)、S (符号, 书写形式)、 G (语法,
排列方法)、P (呈现, 具体记法)。展开序列:

\begin{verbatim}
5 → [5] → {5}
\end{verbatim}

5 是外部符号 (S 层); {[}5{]} 是它指向的概念 (C 层); \{5\} 是沿定义链
追溯到公设基的完整展开 (B 层)。

\subsection{2. 两种 token:
直觉定义与结构展开}\label{ux4e24ux79cd-token-ux76f4ux89c9ux5b9aux4e49ux4e0eux7ed3ux6784ux5c55ux5f00}

\textbf{定义 2.1} (结构 token)。结构 token 是其展开序列完整生成的 token:
概念 → 符号 → 语法 → 呈现, 沿定义链追溯至公设基。展开序列是
\textbf{全部中间层}的显式序列。

\textbf{定义 2.2} (直觉定义 token)。直觉定义 token 是直接由概念指向
生成的 token: 直觉从 C 层接入, 符号/语法/呈现层被\textbf{压缩为占位}。
展开序列中, C 层指向保留, S/G/P 层的展开步骤被跳过。

\textbf{展开序列对比}:

\begin{longtable}[]{@{}lll@{}}
\toprule\noalign{}
维度 & 结构 token & 直觉定义 token \\
\midrule\noalign{}
\endhead
\bottomrule\noalign{}
\endlastfoot
起点 & 概念 (C) & 概念 (C, 直接接入) \\
展开 & 沿定义链逐层 (C→S→G→P→B) & C 层指向 + 表达层占位 \\
中间层 & 全部显式 & 压缩/跳过 \\
终点 & 公设基 (B) & 概念指向 (C) \\
验证 & 展开即验证 & 需外部展开验证 (Lean) \\
\end{longtable}

\textbf{观察 2.3} (本 session): 直觉链 ``复数轴是高维结构中 -1 的投影''
即直觉定义 token --- C 层指向清晰 (投影概念), S/G/P 层 (具体代数 实现)
由 Lean 形式化展开 (ComplexAxis 构造), 全部展开后是 KNOWN 结构
(矩阵表示/代数扩张)。

\subsection{3. 直觉越过什么:
表达层的展开}\label{ux76f4ux89c9ux8d8aux8fc7ux4ec0ux4e48-ux8868ux8fbeux5c42ux7684ux5c55ux5f00}

\textbf{命题 3.1}
(直觉越过的内容)。直觉越过的是\textbf{表达层的展开细节}: - S 层 (符号):
用哪个记号 (a+bJ 还是矩阵) --- 直觉不关心; - G 层 (排列): 参数怎么排
(三元运算槽位) --- 直觉不指定; - P 层 (呈现): 具体记法/优先级 ---
直觉不选择; - \textbf{未越过的}: C 层 (概念指向) --- 直觉从 C 接入,
概念是直觉的 落点。

\textbf{命题 3.2} (越过的代价)。跳过展开 ≠ 免除展开: 直觉的 C 层指向
可能是错 (如本 session 中 ``8 个整点 = 8 个素数''、``发散值 = 单位根''、
``线外无零点很简单'' 等被 Lean 展开证伪的直觉)。跳过展开延后了验证,
验证必须由结构展开 (Lean) 完成。

\textbf{观察 3.3} (语言学类比)。这与生成语法 (Chomsky) 的深层/表层结构
类比: 直觉接入深层结构 (概念/语义), 展开生成表层结构 (符号/语法/
呈现)。语言的直觉理解 = 深层结构直接语义化; 形式化的展开 = 从深层
到表层的生成。token 体系使这一生成可显式化、可验证。

\subsection{4. token
体系为何能产生可解释的结果}\label{token-ux4f53ux7cfbux4e3aux4f55ux80fdux4ea7ux751fux53efux89e3ux91caux7684ux7ed3ux679c}

\textbf{命题 4.1} (可解释性的来源)。token 体系产生可解释结果的机制: 1.
\textbf{层分离}: B/C/S/G/P 五层独立维护 --- 概念 (C) 不依赖具体符号 (S)
或记法 (P), 每个结果的展开序列可追溯至公设基 (B), 因此
``为什么是这个结果''可逐层回答; 2. \textbf{展开可验证}: 每个 token
沿定义链的展开是确定的 (5 → {[}5{]} → \{5\}), 概念对比基于 C 层
(不受记法干扰) --- 结果的 推导链是显式的, 可检查; 3.
\textbf{直觉与展开的解耦}: 直觉在 C 层接入, 展开在 S/G/P 层执行 ---
直觉声称与展开结果可分离比较: 直觉指了概念, 展开给了表达, 两者一致 ⟹
结果可解释。

\textbf{命题 4.2} (为什么直觉命中 KNOWN)。直觉的 C 层接入之所以稳定,
是因为概念层 (数学结构) 是\textbf{稳定对象}: 同一概念有无数表达 (投影 =
矩阵表示 = 代数扩张), 直觉指向概念, 展开选择表达。因此 直觉链命中概念 =
命中该概念的所有表达中的已知结构 (KNOWN) --- 这
解释了''直觉结果为什么是可解释的已知结果'', 而非声称直觉''正确''。

\textbf{观察 4.3} (本 session 证据): C011-C025 全部为 KNOWN 重述 ---
直觉 指向的概念 (投影、两平方和、圆反演、函数方程对称) 都是经典结构的 C
层, 表达层 (ComplexAxis/高斯整数/ℂ) 由 Lean 展开。直觉没有产生 新概念,
但准确选择了概念 --- 结果可解释 (展开链完整), 且与命题 4.2 一致。注意:
直觉也可能选错概念 (观察 3.2), 可解释性不保证 正确性。

\textbf{命题 4.4} (普遍联系与捷径)。概念层 (C) 不是孤立 token 的集合,
而是\textbf{普遍联系的网络}: 概念间的连接 (定义、同构、类比、反演、对偶)
构成图。直觉是图中的\textbf{捷径发现}: 从当前概念直接跳到目标概念, 不经
中间节点的展开。token 体系的独特能力是\textbf{直接定义捷径}: 把直觉
发现的捷径固化为可引用的定义 (新 token 或显式连接), 使捷径成为
网络的结构边, 而非一次性跳跃。

\textbf{观察 4.5} (本 session 的捷径定义): 直觉发现的捷径被逐个固化为
定义/定理: ``复数轴 = 投影'' (C011, 连接高维代数与投影)、``素数 =
圆上整点'' (C015, 连接数论与几何)、``临界线 = 圆'' (C019, 连接函数
方程与反演)、``1/2 = 对称中心'' (C013/C018)。每条捷径都是概念网络的
一条新边, 可引用、可验证、可再展开 --- 这是''直接定义捷径''的实例:
直觉提供捷径 (发现连接), token 体系提供定义 (固化连接)。

\textbf{命题 4.6} (可解释性与捷径)。捷径的可解释性: 每条定义的捷径 (边)
仍保留两端点的完整展开链 (5 → {[}5{]} → \{5\}) --- 捷径是''可解释的
压缩'': 引用时一步直达, 检查时仍可展开验证。这解释了为什么直觉
快速路径的结果可解释: 捷径本身有完整的两端展开, 只是跳过了中间 路径。

\subsection{5. 展开序列的压缩与 token
经济}\label{ux5c55ux5f00ux5e8fux5217ux7684ux538bux7f29ux4e0e-token-ux7ecfux6d4e}

\textbf{观察 5.1} (token 经济)。本 session: 700k token, 99.2\% 缓存命中,
实际新增 \textless{} 1\%。解耦视角的解释: - 直觉定义 token 的展开 (C
层接入) 在\textbf{对话层面}被压缩 --- 概念 直接说出, 展开延后到 Lean; -
缓存命中 (99.2\%) 对应展开序列的重复部分 (同一概念的表达层重发); -
睡眠后的直觉 compact: 时间间隔整理的是 C 层指向 (概念聚焦),
表达层展开被重新生成 --- 对应 §2 的''直觉从 C 接入''。

\textbf{观察 5.2} (展开的延后成本)。延后展开的成本: 直觉错误只有在展开
时暴露 (本 session 多次: 8 点 vs 4 点、共轭对、发散值等修正)。这是
直觉快速路径的固有代价: 跳过表达层展开的验证, 换取即时性。

\subsection{6. 语言学定位}\label{ux8bedux8a00ux5b66ux5b9aux4f4d}

本文框架的语言学定位 (理论提案, 非已证): - \textbf{语义三角} (Odgen \&
Richards, 1923): 符号 (S) ↔ 概念 (C) ↔ 指称 (数学对象)。token
体系的五层是语义三角的细化: S (符号), C (概念), 指称 = 数学对象 (B
层公设基的语义); - \textbf{深层/表层结构} (Chomsky, 1965): 直觉接入深层
(C), 展开生成 表层 (S/G/P); - \textbf{直觉即语义直取} (本提案):
直觉是从概念到指称的直接路径, 不经 表层生成 --- 这正是 token 体系允许的
(C 层独立于表达层)。

\textbf{命题 6.1} (语言学的预测)。若本框架成立, 则: 任何直觉正确指向
概念的系统 (数学直觉、语言直觉), 其错误都发生在表达层展开时
(符号/语法选择), 而非概念层 --- 与观察 3.2 (直觉错误在 Lean 展开时 暴露)
一致。

\subsection{7. 符号定义捷径的观测:
一词多义、多词一义与最快路径}\label{ux7b26ux53f7ux5b9aux4e49ux6377ux5f84ux7684ux89c2ux6d4b-ux4e00ux8bcdux591aux4e49ux591aux8bcdux4e00ux4e49ux4e0eux6700ux5febux8defux5f84}

\textbf{观测是方法}: 符号定义捷径的特点须经观测获得 --- 对比结构 token
与直觉 token 的展开序列, 观察语义与符号的映射过程。

\subsubsection{7.1
语言学的符号定义捷径}\label{ux8bedux8a00ux5b66ux7684ux7b26ux53f7ux5b9aux4e49ux6377ux5f84}

语言学的三个经典现象, 在 token 体系下都是概念网络 (C) 与符号层 (S)
之间的多对多映射:

\begin{longtable}[]{@{}
  >{\raggedright\arraybackslash}p{(\linewidth - 4\tabcolsep) * \real{0.3333}}
  >{\raggedright\arraybackslash}p{(\linewidth - 4\tabcolsep) * \real{0.3333}}
  >{\raggedright\arraybackslash}p{(\linewidth - 4\tabcolsep) * \real{0.3333}}@{}}
\toprule\noalign{}
\begin{minipage}[b]{\linewidth}\raggedright
现象
\end{minipage} & \begin{minipage}[b]{\linewidth}\raggedright
映射
\end{minipage} & \begin{minipage}[b]{\linewidth}\raggedright
token 序列表现
\end{minipage} \\
\midrule\noalign{}
\endhead
\bottomrule\noalign{}
\endlastfoot
一词多义 & 一 S → 多 C & 同一符号 token 展开到不同概念, 由上下文 (前后
token) 决定分支 \\
多词一义 (同义词) & 多 S → 一 C & 不同符号 token 的展开 (\{5\} 追溯)
收敛到同一概念 \\
一义多词 & 一 C → 多 S & 同一概念有多个表达 token (投影 = 矩阵 =
扩张) \\
\end{longtable}

\textbf{观测要点}: 一词多义 = 展开序列的分支点 (同一 S 的 \{5\} 有多个
候选, 上下文收敛); 多词一义 = 展开序列的收敛点 (不同 S 的 \{5\} 汇合);
一义多词 = 表达层的可替换性 (概念稳定, 表达多样)。

\subsubsection{7.2 例: 3+2+7+8
的路径对比}\label{ux4f8b-3278-ux7684ux8defux5f84ux5bf9ux6bd4}

同一目标 (结果 20), 三条 token 序列:

\begin{verbatim}
符号序列: 3, +, 2, +, 7, +, 8

路径 A (顺序展开, 结构): 3+2=5 → 5+7=12 → 12+8=20
  中间节点: {5, 12}

路径 B (结合重排, 结构): (3+7)=10, (2+8)=10 → 10+10=20
  中间节点: {10} (两次出现)

路径 C (直觉, 直接): 28 凑对、37 凑对 → 20
  中间节点: {10 隐式} — 10 作为工作记忆的隐式中间, 不经显式展开
\end{verbatim}

\textbf{观测结果}: 1. \textbf{路径非唯一}: 同一语义 (20) 有多条 token
序列, 中间节点不同 (\{5,12\} vs \{10\} vs 隐式); 2.
\textbf{不必须走近路}: 结构路径 (A/B) 与直觉路径 (C) 都是合法的 ---
直觉选了''最少显式中间节点''的路径 (C), 但 A/B 同样可达; 3.
\textbf{直觉路径的中间节点隐式化}: C 中 10 存在但未展开 (10 分裂成 两个
1 的位值过程被跳过) --- 直觉压缩的是中间概念的显式展开, 不是概念本身。

\subsubsection{7.3 语义理解 =
找到最快路径的数学问题}\label{ux8bedux4e49ux7406ux89e3-ux627eux5230ux6700ux5febux8defux5f84ux7684ux6570ux5b66ux95eeux9898}

\textbf{命题 7.1} (语义理解的数学化)。语义理解是如下最优化问题:
给定符号-逻辑映射 (概念网络 G = (V,E), V = 概念, E = 连接), 输入
符号序列, 求到目标语义的\textbf{最快路径} --- 其中''快''由展开成本定义
(显式中间节点数、工作记忆占用、展开步数)。直觉 = 已学习的最短 路径
(捷径); 结构 = 完整路径; 两者都是 G 中的合法路径, 直觉路径
是优化后的结果。

\textbf{命题 7.2} (可解释性与路径选择)。路径的选择 (A/B/C) 是可解释 的:
每条路径的中间节点序列是可观察的 (token 序列对比), 直觉路径 (C)
的隐式中间节点可通过''追问展开'' (让直觉者显式展开 10) 恢复 --- 这对应
token 体系的展开可追溯性。

\textbf{观测预测} (可检验): 若语义理解 = 最快路径, 则: 1.
同一符号序列存在多条 token 序列, 且可被受试者稳定选择 (个人 偏好路径);
2. 直觉路径的中间节点被隐式化 (工作记忆少), 但可被显式展开恢复; 3.
一词多义的分支选择 = 上下文中的最快路径选择 (哪个候选概念
以最少展开到达当前语境); 4. 多词一义的收敛 =
不同符号路径在概念网络中的汇合点 --- 汇合点 即语义核心 (概念层稳定)。

\subsubsection{7.4 双体系训练: 收敛性设计
(工程推论)}\label{ux53ccux4f53ux7cfbux8badux7ec3-ux6536ux655bux6027ux8bbeux8ba1-ux5de5ux7a0bux63a8ux8bba}

若直觉体系与结构体系共享权重训练, 收敛性设计原则: 1. \textbf{目标分解
(双监督)}: 中间层只由结构路径 (长展开序列) 监督, 末端由直觉路径 (短序列)
与结构路径共同监督 --- 避免两路径对 中间层梯度冲突
(直觉路径的隐式中间节点不参与中间层监督); 2. \textbf{课程顺序}: 先结构
(长展开, 信息完整) 后直觉 (压缩捷径) ---
直觉序列作为已收敛结构表示的捷径蒸馏; 3. \textbf{收敛性}: 在上述设计下,
损失有下界且梯度有界 (标准监督), 收敛可行;
若两路径对同一输入在中间层都强监督, 大概率震荡。

\subsection{8. 词表、形式化强度与信息含量
(猜想)}\label{ux8bcdux8868ux5f62ux5f0fux5316ux5f3aux5ea6ux4e0eux4fe1ux606fux542bux91cf-ux731cux60f3}

\textbf{猜想 8.1} (词表 = 表达空间; 形式化强度 ↔ 紧致性 ↔ 信息含量)。

词表 (符号层 S) 是表达的空间范围。对一种语言: - \textbf{形式化越强,
表达空间越紧致}: 强形式化 (数学记号/Lean) 中 符号-概念映射近乎单射
(一义一词), 展开序列确定 (5 → {[}5{]} → \{5\} 唯一), 表达空间 =
概念空间的紧致表示; - \textbf{形式化越弱, 自由度越多, 空间越大}:
弱形式化 (自然语言) 中 符号-概念映射多对多 (一词多义/多词一义),
展开序列依赖上下文, 表达空间大、自由度多; - \textbf{投影丢失结构的代价}:
弱形式化的符号是强形式化概念的''投影'' (丢失词性/精确语义/语法角色), 单
token 无法确定其概念 --- 需 上下文展开 --- \textbf{单 token
承载的信息含量随形式化减弱而不足}。

\textbf{论证 (信息论)}。单 token 的信息含量 H(token) = -log P(token
\textbar{} 上下文)。强形式化: token 条件分布接近退化 (确定性), 条件熵低,
单 token 结构信息高; 弱形式化: token 上下文依赖强, 边际信息低。

\textbf{例}: - Lean/数学: ``∀'' --- 一个符号, 一个概念, 展开确定 ---
紧致, 高信息; - 自然语言: ``run'' --- 一词多义 (跑步/运行/经营),
上下文决定 --- 松, 低信息; - 3+2+7+8 直觉序列: 压缩 (单 token ``20''
高信息), 但依赖''凑对'' 的上下文知识 --- 介于两者之间。

\textbf{预测 (可检验)}: 对同一语义, 强形式化语言的 token 序列更短、 单
token 展开更确定、信息含量更高; 弱形式化语言的 token 序列更长、 单 token
展开更依赖上下文。可用 token 展开序列的确定性 (条件熵)
量化''形式化强度''。

\textbf{声明}: 本猜想的数学内核 (紧致性/信息含量与符号-概念映射确定性
的关系) 可形式化; 猜想本身是理论提案, 非已证定理。

\subsubsection{8.1 n 基点空间: 语言符号系统的 Lean 模型
(实验)}\label{n-ux57faux70b9ux7a7aux95f4-ux8bedux8a00ux7b26ux53f7ux7cfbux7edfux7684-lean-ux6a21ux578b-ux5b9eux9a8c}

\textbf{定义 8.2} (n 基点空间)。语言符号系统建模为
\texttt{BasepointSpace\ n\ :=\ \{\ basepoints\ :\ Fin\ n\ →\ ComplexAxis\ \}}
--- n 个基点 (直觉定义的锚点) + 投影 (符号层)。实验观测 (Lean 形式化,
全量 build 通过):

\begin{itemize}
\tightlist
\item
  \textbf{松散性定理} (\texttt{pureImag\_collapse}): 纯虚基点 (实部 = 0)
  的投影 全部相同 --- 投影视角下 n 个基点塌缩为一点。直觉越多
  (基点越多), 符号层 (投影) 的区分度越低, 结构越松散;
\item
  \textbf{歧义定理} (\texttt{basepoint\_ambiguity}):
  任意两个不同纯虚基点投影 相同 --- 从投影值无法还原基点,
  符号层信息不足;
\item
  \textbf{不可还原} (\texttt{kernelEquivReal}): 投影 0 的基点族与 ℝ 等势
  (不可数) --- 投影视角的信息不足以区分基点族。
\end{itemize}

\textbf{实验结论} (观测): 在 n 基点模型中, 直觉锚点数量 n 与符号层的
区分度负相关: n 越大, 投影塌缩越严重, 单 token (投影值) 承载的
基点信息越不足 --- 与猜想 8.1 (形式化弱 → 单 token 信息不足) 一致。

\subsubsection{8.2 塌缩投影工具: 任意结构收敛为一点
(规律提取工具)}\label{ux584cux7f29ux6295ux5f71ux5de5ux5177-ux4efbux610fux7ed3ux6784ux6536ux655bux4e3aux4e00ux70b9-ux89c4ux5f8bux63d0ux53d6ux5de5ux5177}

为提取''承载信息与 token 数量/维度''的规律, 需要工具在投影时
\textbf{任意收敛结构} --- 把任意线段、任意多边形、任意高维结构投影为
一个点。Lean 形式化 (全量 build 通过):

\textbf{定义 8.3} (塌缩投影)。对任意载体 S 与指定子结构 K: -
\texttt{collapseRel\ S\ K}: 塌缩等价关系 --- K 中任意两点等价 (K 塌缩),
核外保持恒等; -
\texttt{collapseSpace\ S\ K\ :=\ Quot\ (collapseRel\ S\ K)}: 商空间; -
\texttt{collapseProj\ S\ K\ :\ S\ →\ collapseSpace\ S\ K}: 塌缩投影。

\textbf{定理 8.4} (\texttt{collapse\_kernel}): K
中任意两点在塌缩投影下同像 --- \textbf{任意结构可收敛为一点}。实例:
\texttt{proj} (核 = J 方向) 是 \texttt{collapseProj} 的特例;
任意线段/多边形/高维结构取 K = 全空间即 收敛为单点。

\textbf{规律 (信息与维度)}: 塌缩后核内信息丢失 (核内点不可区分);
商空间可分辨点数 = 核外点数 + 1 (核塌缩为一点)。对应 token 侧: token
维度空间 = 商空间 (投影后结构); 承载信息量 = 核外点数 (tok 数量维度);
直觉 (核) 越大 → 商空间维度越低 → 单 token 信息 越不足 ---
塌缩工具是提取该规律的通用机制, 任意结构可逐级收敛。

\subsection{9. 结论}\label{ux7ed3ux8bba}

token 体系能产生可解释结果的机制: 概念层 (C) 与表达层 (S/G/P) 可 分离,
展开可追溯至公设基 --- 每个结果都有完整的解释链 (为何是它、
从哪来、与什么一致)。直觉快速路径的机制: 直觉从 C 层直接接入,
越过表达层的展开细节 (符号/语法/呈现), 但\textbf{不越过概念本身};
越过的代价是验证延后 (必须由结构展开补验)。本 session 的经验证据 支持:
直觉链命中 KNOWN 概念 (C 层), 表达由 Lean 展开, 直觉错误在 展开时暴露,
睡眠后 C 层指向被整理 (compact)。\textbf{声明边界}: 本文
解释的是''可解释性从何而来'', 不声称 token 体系''正确'' --- 直觉与
展开的一致性是可解释的, 正确性仍是待验证的性质。\textbf{核心机制总结}:
普遍联系 (概念网络) 是捷径的土壤, 直觉是捷径发现者, token 体系是
捷径定义器 (把发现固化为可引用的边) --- 三者结合使直觉快速路径既快
又可解释。

\subsection{附录: 证据来源}\label{ux9644ux5f55-ux8bc1ux636eux6765ux6e90}

\begin{itemize}
\tightlist
\item
  session 2026-08-12: 黎曼方向直觉形式化 (C011-C025, 全部 KNOWN, 无
  sorry, 全量 build 通过)
\item
  token 审计: 700k token, 99.2\% 缓存命中 (session\_token\_audit.py)
\item
  直觉修正案例: 8 点 vs 4 点 (C017 唯一性), 共轭对 (C020/C023), 发散值 =
  单位根 (证伪), 线外无零点 (猜想, 未证)
\end{itemize}

\end{document}
